% Generated draft: requires mathematical and source review.
\documentclass[11pt]{article}
\usepackage[margin=25mm]{geometry}
\usepackage{fontspec}
\setmainfont{Latin Modern Roman}
\usepackage{amsmath,amssymb,mathtools,unicode-math}
\setmathfont{Latin Modern Math}
\usepackage{enumitem,needspace,xurl,hyperref}
\setlength{\parindent}{0pt}
\setlength{\parskip}{5pt}
\setlength{\emergencystretch}{3em}
\newcommand{\N}{\mathbb{N}}
\newcommand{\Z}{\mathbb{Z}}
\newcommand{\Q}{\mathbb{Q}}
\newcommand{\R}{\mathbb{R}}
\newcommand{\C}{\mathbb{C}}
\newcommand{\F}{\mathbb{F}}
\newcommand{\A}{\mathbb{A}}
\newcommand{\PP}{\mathbb{P}}
\newcommand{\E}{\mathbb{E}}
\newcommand{\eps}{\varepsilon}
\newcommand{\dd}{\,\mathrm{d}}
\begin{document}
\section*{UnsolvedMath ID: 2100402}
\textbf{Source problem code:} catalogue entry
\subsection{Definitions and mathematical statement}
The source asks for the algebraic Birkhoff conjecture for outer billiards on a non-Euclidean surface of constant curvature. Thus the curvature is nonzero; the supplied assessment specifically discusses spherical and hyperbolic models.

For geometric context, an outer-billiard step uses a tangent geodesic to the boundary curve $\gamma$ and reflects an exterior point $z$ through the tangency point $m$, so that $m$ is the geodesic midpoint of the step. Choosing tangent branches, orientation and an appropriate exterior domain is necessary to obtain a map $T_\gamma$. On a sphere these choices also involve global hemisphere/domain issues. None of these data is specified in the isolated record.

In a common algebraic-integrability formulation, a nonconstant function $P$ from a specified polynomial class satisfies
\[
 P(T_\gamma(z))=P(z)
\]
on the map's domain, and a rigidity theorem is sought for the boundary curve. However, the source does not define the polynomial class, the ambient model/coordinates, regularity and convexity of $\gamma$, the domain of the invariant, or the exact geometric conclusion of its named conjecture. Polynomiality in ambient coordinates and other possible algebraic integral notions must not be conflated. The record also does not say whether the surface is the simply connected sphere/hyperbolic plane or a constant-curvature quotient.

Consequently the full hypothesis and conclusion of the requested rigidity assertion must be recovered from the source context. No Euclidean theorem, inner-billiard hypothesis, or guessed non-Euclidean rigidity target is silently substituted for it.
\subsection{Short English statement}
Prove the polynomial/algebraic rigidity conjecture for outer billiards in non-Euclidean constant curvature. The supplied assessment is explicitly open/unverified: it reports no verified resolving citation for the spherical/hyperbolic outer-billiard version. It contrasts this with reported planar inner-billiard results of Glutsyuk and Bialy--Mironov and a Euclidean outer-billiard result attributed to Bialy--Mironov, arXiv:2306.12494. Its unsuccessful verification in that session does not establish that no non-Euclidean result exists. The original target remains distinct from those reported Euclidean or inner-billiard cases.
\subsection{Sources}
\begin{itemize}
\item UnsolvedMath permanent ID 2100402, source code catalogue entry, ``Open Problems in Integrable Systems --- Around the Birkhoff conjecture.'' Registry curation and metadata: Ulam AI and UnsolvedMath Contributors, CC BY 4.0; updated 2026-08-06.
\item \emph{Open Problems, Questions, and Challenges in Finite-Dimensional Integrable Systems} (2018), Problem 4.2: \url{https://arxiv.org/abs/1804.03737}. Supplied source-TeX extraction accessed 2026-07-29. The metadata says the source presents an open item; status/release and rights review are NEEDS\_REVIEW.
\item Bialy--Mironov, arXiv:2306.12494, \url{https://arxiv.org/abs/2306.12494}, as cited for the Euclidean outer-billiard comparison; the summary points to source code AMR-020-0401. No article title or resolving non-Euclidean reference is supplied.
\item Glutsyuk and Bialy--Mironov are named for planar inner-billiard algebraic results without complete bibliographic entries. No missing titles, dates, or locators are inferred.
\end{itemize}
\end{document}
